\scanpage{391}
\mspage{\deleted{\unsure{116}}115}
% archive sheet number: 390
qu'il suffit de prouver qu'il en est ainsi si $X$ a un objet initial (axiome \inserted{$\mathrm{Loc}(2^*)$} dual de l'axiome $\mathrm{Loc}(2)$). A ce moment, on saura que $W$ est un localiseur fondamental, et de plus il satisfait $\mathrm{Loc}(3^{\unsure{*}}\mathrm{bis})$ donc aussi $\mathrm{Loc}(3\,\mathrm{bis}\unsure{)}$ (forme forte de l'axiome de localisation).

Donc \uline{\unsure{DPS|IPS|APS}} que \deleted{$\unsure{S}$} si $S$ a un objet \unsure{initial}, alors $S$ est $W$-asphérique, i.e. \deleted{le foncteur $p:$} $S\unsure{\relbar|\longrightarrow}e$ est dans $W$, i.e $p^{\unsure{\wedge|*}}:D(e)\relbar D(S)$ est pl. fidèle.
Il revient au \ensuremath{=} de dire que la flèche d'adjonction
\[
p_!p^*\unsure{\xi|\zeta}\longrightarrow\xi
\]
est un iso, pour \unsure{tt} $\xi$ dans $D(e)$. Or le calcul de $p_!(y)$ (pour $y\in\mathrm{Ob}\,D(S)$\inserted{ i.e. $y\in\mathrm{Ob}\,M(S)$})\unsure{,} se fait ainsi : on choisit une flèche
\[
y'\xrightarrow{\Sigma_S}y^{\deleted{\unsure{\prime|*}}}\,,
\]
avec $y'$ ``cofibrant'' (pour la structure à cofibrations pertinente sur $M(S)$), et on
